\chapter{Biology and Technology}\label{ch:g9:biotech}

\omterm{def:g6:method:science}{Biology} and technology meet in medicine, \omterm{def:g2:food-energy:food}{food}, forensics and
\omterm{def:g7:impact:conserve}{conservation}. This chapter introduces tools and ethical questions at a
first-look level --- not a lab manual for advanced genetic engineering.

\section{What counts as biotechnology?}

\begin{definition}[Biotechnology]\label{def:g9:biotech:def}
\emph{Biotechnology}\index{biotechnology} is the use of living \omterm{def:g4:ecosystems:levels}{organisms},
\omterm{def:g7:microscope:cell}{cells} or biological molecules to make products or solve problems --- from
traditional yoghurt and bread to modern \omterm{def:g9:dna:dna}{DNA} tools.
\end{definition}

\begin{example}[Traditional biotech]\label{ex:g9:biotech:trad}
Yeast ferments dough and brews; \omterm{def:g4:microbes:def}{bacteria} ferment milk into yoghurt and
cheese. Humans selected \omterm{def:g4:microbes:def}{microbes} long before \omterm{def:g9:dna:gene}{genes} were named.
\end{example}

\begin{omfigure}
\begin{tikzpicture}[font=\footnotesize,
  b/.style={draw, rounded corners, align=center, minimum width=2.3cm,
    minimum height=0.85cm}]
  \node[b, fill=omDef!15] (t) at (0,0) {traditional\\fermentation};
  \node[b, fill=omMeth!15] (m) at (3.5,0) {medical\\(insulin, vaccines)};
  \node[b, fill=omProp!15] (a) at (7,0) {agriculture\\\& environment};
  \node at (3.5,1.2) {biotechnology spans many fields};
\end{tikzpicture}

{\small Broad map of applications.}
\end{omfigure}

\section{DNA technologies (ideas)}

\begin{definition}[Genetic testing and profiling]\label{def:g9:biotech:test}
\emph{DNA profiling}\index{DNA profiling} compares patterns in \omterm{def:g9:dna:dna}{DNA} to
match samples (forensics, kinship). Genetic tests can detect \omterm{def:g7:heredity:allele}{alleles}
linked to some diseases --- with counselling and privacy \omterm{def:g1:caring:needs}{needs}.
\end{definition}

\begin{definition}[Genetic modification idea]\label{def:g9:biotech:gm}
\emph{Genetic modification}\index{genetic modification} changes an
\omterm{def:g4:ecosystems:levels}{organism}'s \omterm{def:g9:dna:dna}{DNA} deliberately (for example inserting a \omterm{def:g9:dna:gene}{gene} so \omterm{def:g4:microbes:def}{bacteria} make
human insulin). Regulations and risk assessments differ by country and
use.
\end{definition}

\begin{example}[Insulin]\label{ex:g9:biotech:insulin}
\omterm{def:g4:microbes:def}{Bacteria} engineered with the human insulin \omterm{def:g9:dna:gene}{gene} produce insulin for
diabetes treatment --- replacing limited \omterm{def:g1:plants-animals:animal}{animal} extracts.
\end{example}

\begin{definition}[Cloning (everyday meanings)]\label{def:g9:biotech:clone}
In biotech talk, \emph{cloning}\index{cloning} can mean copying \omterm{def:g9:dna:dna}{DNA}
fragments in \omterm{def:g4:microbes:def}{bacteria}, or making a genetically similar \omterm{def:g4:ecosystems:levels}{organism}. Context
matters; \omterm{def:g6:method:science}{science} fiction is not the lab notebook.
\end{definition}

\begin{omfigure}
\begin{tikzpicture}[font=\footnotesize]
  \draw[thick, omThm, fill=omThm!12] (0,0) rectangle (1.8,1.2);
  \node[align=center] at (0.9,0.6) {gene of\\interest};
  \draw[-Stealth, thick] (2.0,0.6) -- (3.0,0.6);
  \draw[thick, omDef, fill=omDef!15] (3.2,0.1) ellipse (1.1 and 0.7);
  \node at (3.2,0.6) {microbe};
  \draw[-Stealth, thick] (4.5,0.6) -- (5.5,0.6);
  \draw[thick, omProp, fill=omProp!20] (5.7,0.2) rectangle (7.5,1.0);
  \node[align=center] at (6.6,0.6) {useful\\protein};
\end{tikzpicture}

{\small Idea: insert instructions $\to$ \omterm{def:g7:microscope:cell}{cell} factory $\to$ product
(schematic).}
\end{omfigure}

\section{Selective breeding vs modern tools}

\begin{proposition}[Same goal, different speed]\label{prop:g9:biotech:breed}
\emph{Selective breeding}\index{selective breeding} chooses parents with
desired \omterm{def:g7:heredity:trait}{traits} over generations. Molecular methods can move specific \omterm{def:g9:dna:gene}{genes}
faster, but still interact with whole \omterm{def:g4:ecosystems:levels}{organisms} and \omterm{def:g7:eco-energy:ecosystem}{ecosystems}.
\end{proposition}
\admitted

\begin{method}[Ethical checklist]\label{met:g9:biotech:ethics}
\begin{enumerate}
  \item Who benefits and who bears risk?
  \item Is consent and privacy respected (human data)?
  \item What are ecological effects if \omterm{def:g4:ecosystems:levels}{organisms} are released?
  \item Are claims supported by evidence, not only marketing?
\end{enumerate}
\end{method}

\section{Exercises}

\begin{exercise}[$\star$]\label{exo:g9:biotech:1}
Define \omterm{def:g9:biotech:def}{biotechnology} and give one traditional example.
\end{exercise}

\begin{exercise}[$\star$]\label{exo:g9:biotech:2}
How do \omterm{def:g4:microbes:def}{bacteria} help make human insulin today (idea level)?
\end{exercise}

\begin{exercise}[$\star$]\label{exo:g9:biotech:3}
What is \omterm{def:g9:biotech:test}{DNA profiling} used for?
\end{exercise}

\begin{exercise}[$\star$]\label{exo:g9:biotech:4}
Distinguish \omterm{prop:g9:biotech:breed}{selective breeding} from \omterm{def:g9:biotech:gm}{genetic modification} in one sentence
each.
\end{exercise}

\begin{exercise}[$\star$]\label{exo:g9:biotech:5}
Why does genetic privacy matter?
\end{exercise}

